\section{Introduction}
\label{sec:intro}

Vehicles generate compute tasks that they cannot process onboard and can offload to roadside fog nodes (RSUs); which RSU receives a task is the decision studied here.
A certificate is evidence about identity: it says who a node is.
Whether the node is \emph{currently working correctly} is a different question, and one that observed task outcomes can speak to.
This paper calls the difference between the two the certificate-versus-behaviour gap.

As a published illustration that V2X software can stop working:
% src: A26; CVE_VERIFICATION.md (read 2026-10-15). Wording of A26 only (X13, X20): nothing about ITS-G5, radio range, deployed RSUs, field use, certificate holders or SCMS.
CVE-2026-37554, CVE-2026-43988 and CVE-2026-44905 are published CVE records of denial-of-service vulnerabilities (CVSS 3.1 base 7.5) in the open-source Vanetza V2X implementation ($\le$ 26.02), in each of which a remote unauthenticated network attacker can terminate the process through an uncaught exception (CWE-248): the first in ECC point validation in GeoNetworking packet processing (published 2026-05-01), the second in ASN.1/OER parsing and the third in the cryptographic verification path (both published 2026-05-26)~\cite{cve1,cve2,cve3}.
The records were read from the CVE Program record service on 2026-10-15.
We draw no further conclusion from these records.

We add a behavioural trust layer to the offloading decision.
A trust head scores each RSU from observed task outcomes, and a fixed analytic rule, $\text{score}=\alpha\,\text{trust}-\beta\,\text{latency}-\gamma\,\text{load}$, selects the target among in-range, dwell-feasible RSUs.
% src: A1
Selection is not learned, and the baseline without trust is the same rule with $\alpha=0$.
The trust head is trained self-supervised to predict whether an RSU meets its next task deadline; the injected behaviour class is never a label or a feature.
% src: A2
We compare a node-local head (C) with a graph neural network head (D) that aggregates evidence over the network graph, under the same rule, the same features and the same training objective.
The question the comparison is built to answer is whether message passing over the graph adds anything beyond the node's own evidence.

The evaluation was designed so that a null or negative answer could be reported at full strength.
Each comparison has a rule (metric, floor, window, verdict labels) that was written to the repository, and committed, before the numbers it judges (Sec.~\ref{sec:protocol}).
All results are for one simulated scenario with injected degradation, and the text says so throughout.

\paragraph*{Contributions}
In this simulation, the paper contributes:
% src: A21. The list is A21 and nothing more; it must not be read as implying that the graph helped (X1, X2).
\begin{enumerate}
\item a behavioural trust layer that drives a fixed weighted selection rule, validated in the closed loop against a tuned no-trust baseline (A) and a perfect-knowledge oracle upper bound (B-oracle), in which the node-local head (C) is the strongest;
\item an exactly decomposable explanation of each selection (Layer 1);
\item pre-registered negative and null results on graph structure, with the measured boundary condition of this scenario.
\end{enumerate}

The rest of the paper follows the order in which the questions were asked: related work (Sec.~\ref{sec:related}), the scenario (Sec.~\ref{sec:system}), the method and its baselines (Sec.~\ref{sec:method}), the explanation layers (Sec.~\ref{sec:explanation}), the protocol (Sec.~\ref{sec:protocol}), results (Sec.~\ref{sec:results}), discussion, limitations and conclusion.
